\exnumber{02}
\extitle{UX V - Ex 3 - Jobs to be done}
\exfiles{ex02/PDF, nomme Name\_FirstName\_DeliverableName\_Date\_VersionNumber}
\exauthorize{Ceux nommés dans l'exercice}
\makeheaderfiles

\textbf{Instructions :}

Pratiquons Jobs-to-be-Done pour être plus précis sur la façon dont un site de label musical pourrait être amélioré.

Lisez à propos de Jobs-to-be-Done et de la dimension psychologique derrière cet outil.

Ensuite, trouvez quelqu'un (votre « utilisateur ») autour de vous qui aime beaucoup la musique.

Demandez-leur ce qu'ils feraient sur le site d'un label musical. Demandez-leur pourquoi, pourquoi et pourquoi pour comprendre l'émotion derrière une action, et le but derrière une simple visite sur un site web.

Complétez le template JTBD ci-dessous en incluant le véritable objectif de votre « utilisateur ».

Et rappelez-vous : vous n'êtes pas votre utilisateur !

\textbf{Livrable :}

\begin{itemize}
\item Un Jobs-to-Be-Done complet sur un document PDF
\begin{itemize}
\item « Quand \_\_\_\_\_\_, je veux \_\_\_\_\_ afin de \_\_\_\_\_ »
\end{itemize}
\item Déposez le PDF sur Github, nommé ainsi : « Name\_FirstName\_DeliverableName\_Date\_VersionNumber ». Exemple : « Doe\_John\_MidFiPrototype\_05242024\_V1 ».
\item Ajoutez un titre dans le document.
\item Soignez l'aspect global du document. Il doit être clair, simple et lisible. Vous pouvez regarder canva.com, sans surcharger le dessin de détails.
\end{itemize}

\textbf{Ressources :}

\begin{itemize}
\item Qu'est-ce qu'un job to be done ? \url{https://jobs-to-be-done.com/what-is-jobs-to-be-done-fea59c8e39eb}
\item Framework Jobs-to-be-done \url{https://www.productplan.com/glossary/jobs-to-be-done-framework/}
\item https://jobs-to-be-done.com/ \url{https://jobs-to-be-done.com/}
\item « Quand nous achetons un produit, nous le « recrutons » essentiellement pour nous aider à accomplir un travail. S'il fait bien le travail, la prochaine fois que nous sommes confrontés au même travail, nous avons tendance à recruter ce produit à nouveau. S'il fait un travail médiocre, nous le « licenciions ». » Clayton M Christensen « Know your customers: JtBD »
\end{itemize}

\textbf{Citation :}

\begin{itemize}
\item « Quand nous identifions les user stories tôt dans le processus de design, nous les laissons dicter les décisions de conception. Plutôt que de laisser nos idées préconçues sur ce que le produit devrait être ou devrait ressembler le faire. » UX Booth
\item Documenter le quoi et le pourquoi de chaque élément favorise l'organisation et rend la transmission à l'équipe de développement bien plus fluide.' UX Booth
\end{itemize}

